% Page 003

\seite{3}

\abschnitt{1. Lage}
\pstart
\trenner


\pend
\begin{verse}
„Wer kennt die Völker, nennt die Namen,\\
Die gastlich hier zusammen kamen."
\end{verse}
\pstart


\margmark{Die Lage des\\Dorfes.}%
    Das Dörfchen Sefferweich liegt unter 50° 3' nördlicher\ob
Breite und 24° 10' \darueber{6° 32'} östlicher Länge im Kreise Bitburg, Bezirk\ob
Trier, ungefähr 8 Km in nördlicher Richtung von der\ob
Kreisstadt Bitburg entfernt, in einem engen ein\ohb
geschlossenen Thälchen, fünfzehn Minuten westlich von\ob
der Trier-Aachener Straße entfernt, in der Bürger\ohb
meisterei Bickendorf, von Wiesen rings umgeben.\ob
Nahe an das Dorf herantretend, zeigt sich ein Wald, \enquote{Brohselt}\ob
\darueber{früher} genannt, von mächtigen Eichen und Buchen bestanden;\ob
von den ersteren sind sehr wenige mehr vorhanden,\ob
\margmark{Neubau der\\Pfarrkirche zu\\Seffern.}%
sie sind zu dem, im Jahre 1851 in Angriff genom\ohb
menen Neubau der Pfarrkirche zu Seffern ver\ohb
wandt worden. Auch die mächtigen Buchen mußten\ob
einem neuen Holzgewächs die Stelle räumen. Der\ob
oben bezeichnete Wald dehnt sich beinahe eine Meile\ob
in nördlicher Richtung, bis nach Balesfeld im\ob
Kreise Prüm aus und schließt sich an den,\ob
später zur Besprechung kommenden Kyllwald\ob
an. An denselben, ungefähr 10 Minuten vom Dorfe\ob
\margmark{Der „Siebenge-\\meinden-Wald".}%
entfernt, beginnt der Neidenwald, eine Lohhecke. An\ob
demselben partizipieren die Gemeinden: Bitburg,\ob
Masholder, Irsch, Fließem, Nattenheim, Malberg
\pend
